\begin{textsample}{NEG Dim 3 – unprofiled\_gpt – Score -2.00 – unprofiled\_gpt/t000315\_gpt.txt}  \label{ex:f3_neg_023}
Yeah, I mean, I ’d already heard that kind of thing even before I came to Poland, this whole narrative about Syrian refugees being rapists.
It was interesting – and depressing – that one of my students this week started going on about “ all these Syrian immigrants coming over here raping women. ” It rang a bell because that exact perception was already floating around in the media before I ever moved here.
What strikes me is how quickly “ Syrian ” became shorthand for “ dangerous ” in people ’s minds, when actually, like you ’re saying, the statistics and the reporting often point somewhere else – North African migrant workers, different contexts, different issues.
But on the ground, in everyday conversations, it just gets boiled down to “ Syrians are rapists, ” and that ’s the version my student has picked up and is repeating as if it ’s \textbf{fact}.
It ties back to what we were both saying about being immigrants ourselves.
I call myself an immigrant quite happily, but I ’m very aware of how loaded that word is when it ’s applied to others, especially people from Syria, for example, who absolutely aren't “ \textbf{okay} ” in the way you and I are treated here.
So when I hear my student talk like that, it really underlines that the attitude was here long before this government – the election just sort of brought it to the surface rather than created it.

% matched lemmas: fact, okay
\end{textsample}
